\chapter{Kursidee und Ablauf}

\section{Worum es geht}
Die SchülerInnen lernen, wie sich Anlageformen unterscheiden, und simulieren selbst, was Unsicherheit über viele Jahre bewirkt. Jede Doppellektion beginnt mit einer Frage, auf die sie noch keine Antwort haben. Danach arbeiten sie selbstständig in vorbereiteten Notebooks. Am Schluss modellieren sie eine eigene reale Entscheidung und zeigen, wo ihr Modell versagt.

Der Kurs ist ein Spezialangebot der Informatik und wird von einer Lehrperson geleitet. Die Fragen stammen aus Wirtschaft und Finanzen, bearbeitet werden sie mit Simulationen in Python. Die nötigen Begriffe aus Wirtschaft und Recht führt der Kurs selbst ein, Vorwissen aus diesem Fach wird nicht vorausgesetzt. Der Kurs gibt keine Anlageempfehlungen.

\section{Umfang}
Der Kurs ist ein Förderangebot im Sinn des Konzepts Begabtenförderung. Er ergänzt den Regelunterricht und umfasst ein Quartal.

\begin{center}
	\small
	\begin{tabular}{@{}p{.3\textwidth}p{.6\textwidth}@{}}
		\midrule
		Präsenz & 8 Doppellektionen zu 90 Minuten, also 16 Lektionen \\
		Eigenarbeit & rund 8 Stunden für das Projekt \\
		Abschluss & Eigene Simulation mit Präsentation \\
		\midrule
	\end{tabular}
\end{center}

\section{Unterrichtsform}
Pro Doppellektion gibt es nur einen kurzen Impuls. Die SchülerInnen arbeiten zu zweit in einem Notebook, das Anleitung, Hilfestellungen und Selbstkontrollen enthält. Wer schneller ist, wählt die Challenges. Die detaillierte Planung steht im Anhang.

\begin{center}
	\small
	\begin{tabular}{@{}lrr@{}}
		\midrule
		\textbf{Phase} & \textbf{Minuten im ganzen Kurs} & \textbf{Anteil} \\
		\midrule
		Einstiegsfrage und Impuls & 105 & 15 \% \\
		Selbstständige Arbeit und Präsentationen & 510 & 71 \% \\
		Austausch in der Gruppe & 105 & 15 \% \\
		\midrule
	\end{tabular}
\end{center}

\section{Material}
Die Notebooks liegen unter \texttt{FF/Risiko\_simulieren/Notebooks}, die Lösungen im Unterordner \texttt{Loesungen}. Die Notebooks 1 bis 7 gehören zu den Doppellektionen 1 bis 7. Notebook 8 ist die Projektvorlage.
